% SPDX-FileCopyrightText: 2026 DX-Lab contributors
% SPDX-License-Identifier: CC-BY-4.0
%
% Báo cáo kỹ thuật DX-Lab v0.2.1. Biên dịch: tectonic -X compile dx-lab-bao-cao-ky-thuat.tex
\documentclass[11pt,a4paper]{article}

\usepackage{fontspec}
\setmainfont{Times New Roman}
\setsansfont{Arial}
\setmonofont{Menlo}[Scale=0.82]

\usepackage[a4paper,left=2.5cm,right=2cm,top=2.2cm,bottom=2.2cm]{geometry}
\usepackage{microtype}
\usepackage{booktabs,tabularx,longtable,array}
\usepackage{amsmath}
\usepackage{xcolor}
\usepackage{graphicx}
\usepackage{enumitem}
\usepackage{fancyhdr}
\usepackage{titlesec}
\usepackage{fancyvrb}
\usepackage{tikz}
\usetikzlibrary{positioning,arrows.meta,shapes.geometric,fit,backgrounds,calc}
\usepackage[hidelinks,pdftitle={DX-Lab v0.2.1: báo cáo kỹ thuật},pdfauthor={HaUI.NPC}]{hyperref}

\definecolor{dxblue}{HTML}{1F4E79}
\definecolor{dxgrey}{HTML}{F2F2F2}
\definecolor{dxred}{HTML}{9C2B2B}
\definecolor{t1}{HTML}{DCE9F5}
\definecolor{t2}{HTML}{E4F1DD}
\definecolor{t3}{HTML}{FBEFD5}

\setlist{itemsep=1pt,topsep=3pt,parsep=0pt}
\setlength{\parindent}{0pt}
\setlength{\parskip}{2.5pt}
\renewcommand{\arraystretch}{1.15}
\setlength{\emergencystretch}{2em}

\titleformat{\section}{\sffamily\Large\bfseries\color{dxblue}}{\thesection}{0.6em}{}
\titleformat{\subsection}{\sffamily\large\bfseries\color{dxblue}}{\thesubsection}{0.6em}{}
\titlespacing*{\section}{0pt}{12pt}{4pt}
\titlespacing*{\subsection}{0pt}{8pt}{2pt}

\pagestyle{fancy}
\fancyhf{}
\renewcommand{\headrulewidth}{0.3pt}
\fancyhead[L]{\sffamily\small DX-Lab v0.2.1 --- báo cáo kỹ thuật}
\fancyhead[R]{\sffamily\small HaUI.NPC · OLP PMNM 2026}
\fancyfoot[C]{\sffamily\small\thepage}

\DeclareUrlCommand\codeX{\urlstyle{tt}\def\UrlBreaks{\do\/\do\_\do\.\do\-\do\,}\def\UrlBigBreaks{\do\/}\def\UrlSpecials{\do\ {\ }}}
\newcommand{\code}{\protect\codeX}
\newcommand{\gd}{\textit{giả định}}
\newcolumntype{L}[1]{>{\raggedright\arraybackslash}p{#1}}
\newcolumntype{Y}{>{\raggedright\arraybackslash}X}
\newcommand{\tbl}[1]{\small #1}

\renewcommand{\contentsname}{Mục lục}
\renewcommand{\figurename}{Hình}
\renewcommand{\tablename}{Bảng}
\newcommand{\cmd}[1]{\texttt{\detokenize{#1}}}
\setcounter{tocdepth}{1}
\begin{document}

%================ BÌA ================
\begin{titlepage}
\centering
\vspace*{1.5cm}
{\sffamily\color{dxblue}\fontsize{34}{40}\selectfont\bfseries DX-Lab\par}
\vspace{0.4cm}
{\sffamily\Large Trạm điều hành số trên nền tảng\\DX-OS Open-Core\par}
\vspace{0.8cm}
{\sffamily\Huge\bfseries Báo cáo kỹ thuật\par}
\vspace{1.2cm}
\rule{0.6\textwidth}{0.6pt}\par\vspace{0.8cm}
{\large
\begin{tabular}{rl}
Đội thi: & \textbf{HaUI.NPC}\\
Cuộc thi: & Olympic Tin học Sinh viên Việt Nam 2026, khối Phần mềm nguồn mở\\
Chủ đề: & DX-OS Open-Core\\
Phiên bản: & \textbf{v0.2.1} (tag \code{v0.2.1}, commit \code{3b06a1e})\\
Ngày phát hành: & 01/11/2026 (giờ Việt Nam)\\
\end{tabular}}\par
\vspace{1.2cm}
\begin{tabular}{rl}
Mã nguồn: & \url{https://github.com/HaUI-NPC/dx-lab}\\
Bản phát hành: & \url{https://github.com/HaUI-NPC/dx-lab/releases/tag/v0.2.1}\\
\end{tabular}\par
\vfill
\begin{minipage}{0.85\textwidth}\small\centering
Giấy phép: mã nguồn Apache-2.0, tài liệu (gồm báo cáo này) CC BY 4.0.\\
Mọi khẳng định dẫn tới tệp trong repo để kiểm lại. Những gì chưa có được nêu riêng ở mục~\ref{sec:lotrinh}.
\end{minipage}
\end{titlepage}

%================ TÓM TẮT + MỤC LỤC ================
\section*{Tóm tắt}
\addcontentsline{toc}{section}{Tóm tắt}
DX-Lab số hóa quy trình xử lý sự cố thiết bị phòng thực hành, từ báo hỏng, định tuyến và xếp mức ưu tiên, kiểm tra,
duyệt chi phí theo hạn mức, sửa chữa tới nghiệm thu. Dự án trả lời yêu cầu ``Open-Core trước, ứng dụng sau'' của đề thi
bằng cách tách hai việc. \textbf{Tầng 1} là nền tảng headless dùng chung (Keycloak, Apache APISIX, PostgreSQL, Core API
FastAPI nhúng SpiffWorkflow) không biết ứng dụng nào đang chạy. \textbf{Tầng 3} là các ứng dụng (\code{lab-repair},
\code{room-booking}) chỉ gồm tệp BPMN, DMN và JSON Schema nạp qua API. Giữa hai tầng, \textbf{Tầng 2} cung cấp cổng web
(portal), máy chủ MCP và CLI \code{dxlab}, đều gọi Core bằng token của chính người dùng. Core kiểm quyền theo năm lớp,
mặc định từ chối, hỗ trợ ủy quyền có thời hạn (TUQ.). Phần trí tuệ nhân tạo gồm bảng luật DMN có mã lý do và tìm kiếm
kết hợp toàn văn với embedding chạy cục bộ; số đo thật: \code{tech_unit} đúng 57/60, mức ưu tiên đúng 53/60, tìm kiếm
top-1 đúng 18/18 trên bộ câu hỏi tự soạn. Bản \code{v0.2.1} đã gắn tag, kèm SBOM và \code{SHA256SUMS}. Các giới hạn
(dữ liệu đánh giá tổng hợp, hạn mức và tổ chức là giả định, chưa có LLM, chưa có màn hình tạo ủy quyền) được nêu rõ.

\tableofcontents
\newpage

%================ 1 ================
\section{Bài toán và phạm vi}
Sinh viên và giảng viên phát hiện thiết bị phòng lab hỏng nhưng thường báo qua tin nhắn hoặc lời nói. Yêu cầu dễ thất
lạc, không ai theo dõi được ai đang xử lý, chi phí sửa chữa không đi qua đúng cấp duyệt, còn tri thức sửa chữa nằm trong
đầu một vài cán bộ kỹ thuật (\code{README.md}, mục ``Bài toán'').

Đề thi yêu cầu \textbf{Open-Core trước, ứng dụng sau}: không dừng ở một ứng dụng web có dashboard hoặc AI. Vì vậy
DX-Lab tách \emph{nền tảng dùng chung} (định danh/SSO, cổng API, dữ liệu, bộ máy quy trình) khỏi \emph{ứng dụng minh họa}
(\code{lab-repair}: sáu bước từ báo hỏng tới nghiệm thu; \code{room-booking}: mượn phòng thực hành). Phạm vi bản
\code{v0.2.1} theo \code{README.md} (mục ``Trạng thái'') và \code{CHANGELOG.md} (mục \code{[0.2.1]}): chạy đầu-cuối trên
một máy bằng Docker Compose với dữ liệu và tài khoản demo.

\section{Kiến trúc Open-Core ba tầng}
\begin{figure}[ht]
\centering
\begin{tikzpicture}[
  font=\sffamily\footnotesize,
  box/.style={draw,rounded corners=2pt,minimum height=0.8cm,align=center,inner sep=3pt},
  tier/.style={draw,rounded corners=4pt,inner sep=7pt},
  arr/.style={-{Stealth[length=2.2mm]},thick,dxblue},
  darr/.style={-{Stealth[length=2mm]},dashed,gray}]
% Tầng 3
\node[box,fill=white,text width=3.6cm] (lr) {\textbf{lab-repair}\\BPMN, 4 DMN, 7 form};
\node[box,fill=white,text width=3.6cm,right=0.6cm of lr] (rb) {\textbf{room-booking}\\BPMN, DMN, form};
\begin{scope}[on background layer]
\node[tier,fill=t3,fit=(lr)(rb),label={[anchor=south west,inner sep=2pt]north west:\bfseries Tầng 3. Ứng dụng}] (T3) {};
\end{scope}
% Tầng 2
\node[box,fill=white,text width=2.4cm,below=1.5cm of lr,xshift=-0.9cm] (por) {\textbf{Portal}\\Prefab UI + BFF};
\node[box,fill=white,text width=2.4cm,right=0.5cm of por] (mcp) {\textbf{MCP server}\\5 công cụ, stdio};
\node[box,fill=white,text width=2.4cm,right=0.5cm of mcp] (cli) {\textbf{CLI dxlab}\\login, validate, publish};
\begin{scope}[on background layer]
\node[tier,fill=t2,fit=(por)(mcp)(cli),label={[anchor=south west,inner sep=2pt]north west:\bfseries Tầng 2. H-P-D-I}] (T2) {};
\end{scope}
% Tầng 1
\node[box,fill=white,text width=2.4cm,below=1.7cm of mcp] (gw) {\textbf{APISIX}\\JWT, rate limit};
\node[box,fill=white,text width=2.4cm,right=0.9cm of gw] (core) {\textbf{Core API}\\FastAPI + SpiffWorkflow};
\node[box,fill=white,text width=2.4cm,left=0.9cm of gw] (kc) {\textbf{Keycloak}\\SSO, OIDC, OTP};
\node[box,fill=white,text width=2.4cm,below=0.7cm of core] (pg) {\textbf{PostgreSQL}\\dữ liệu, FTS};
\begin{scope}[on background layer]
\node[tier,fill=t1,fit=(kc)(gw)(core)(pg),label={[anchor=south west,inner sep=2pt]north west:\bfseries Tầng 1. Open-Core headless}] (T1) {};
\end{scope}
\draw[darr] ($(T3.south)+(3.5,0)$) -- node[right,font=\sffamily\scriptsize]{dxlab publish} ($(T2.north)+(3.5,0)$);
\draw[arr] (T2.south -| gw) -- node[right,font=\sffamily\scriptsize,pos=0.3]{Bearer token người dùng} (gw.north);
\draw[arr] (gw) -- (core);
\draw[arr] (core) -- (pg);
\draw[darr] (kc) -- node[above,font=\sffamily\scriptsize]{JWKS} (gw);
\end{tikzpicture}
\caption{Ba tầng và hướng phụ thuộc (thư mục \code{apps/}, \code{data/}; \code{portal/}, \code{mcp/}, \code{core/}; \code{deploy/}, \code{core/}), rút gọn từ \code{docs/ARCHITECTURE.md} mục 1. Mũi tên liền là lời gọi chạy; nét đứt là nạp định nghĩa hoặc khóa công khai.}
\end{figure}

\subsection{Tầng 1: Open-Core headless}
\begin{center}\tbl
\begin{tabularx}{\textwidth}{@{}L{2.6cm}L{3.6cm}Y@{}}
\toprule
Dịch vụ chung & Thành phần (giấy phép) & Vai trò \\
\midrule
Identity / SSO & Keycloak 26.8.1 (Apache-2.0) & Đăng nhập, nhóm theo cây tổ chức, OTP cho vai trò duyệt, cấp JWT \\
API Gateway & Apache APISIX 3.19.0 (Apache-2.0) & Cửa duy nhất vào Core: kiểm JWT bằng JWKS, giới hạn tốc độ, gắn \code{X-Request-Id} \\
Dữ liệu & PostgreSQL 17.11 (PostgreSQL License) & Phiếu, bước, danh mục, audit, tìm kiếm toàn văn \\
Workflow & SpiffWorkflow nhúng trong Core (LGPL-3.0) & Chạy BPMN 2.0, đánh giá DMN 1.3, lưu trạng thái bền vững \\
\bottomrule
\end{tabularx}
\end{center}
Bốn điểm kiến trúc then chốt (\code{docs/ARCHITECTURE.md} mục 1--3): (1) Core không trả HTML, không biết tên ứng dụng,
không gọi LLM; (2) chỉ Core đọc/ghi PostgreSQL, cổng của Core không bind ra máy chủ, APISIX là điểm vào duy nhất
\code{127.0.0.1:9080}; (3) ranh giới import được giữ bằng khai báo phụ thuộc riêng của từng gói (\code{dxlab-core} không
import \code{dxlab_portal}, \code{prefab_ui}, \code{fastmcp}; \code{dxlab-portal} không import \code{dxlab_core},
\code{SpiffWorkflow}, \code{sqlalchemy}); (4) API công bố bằng OpenAPI tại \code{/api/openapi.json}.

\subsection{Tầng 3: ứng dụng không sửa Core}
\code{apps/lab-repair/} gồm \code{lab-repair.bpmn}, các bảng luật \code{triage.dmn}, \code{routing.dmn},
\code{authority.dmn}, \code{overrun.dmn}, bảy cặp \code{*-schema.json}/\code{*-uischema.json} và bản ghi danh mục.
\code{apps/room-booking/} có \code{room-booking.bpmn}, \code{conflict.dmn} và form. Thư mục \code{apps/} chỉ chứa
\code{.bpmn}, \code{.dmn}, \code{.json}, \code{.md}; validator từ chối gói có script. Ứng dụng thứ hai được nạp qua
\code{/api/definitions} mà không đổi \code{core/} hay \code{portal/} (\code{apps/room-booking/README.md}). Đổi ứng dụng thì
Core không đổi, trừ khi cần một năng lực nền tảng mới dùng chung.

\subsection{Luồng gọi dịch vụ}
Mọi thao tác nghiệp vụ là một lời gọi \code{/api/...} qua APISIX. Ví dụ báo hỏng từ portal
(\code{docs/ARCHITECTURE.md} mục 3.1--3.5):
\begin{figure}[ht]
\centering
\begin{tikzpicture}[font=\sffamily\scriptsize,
  life/.style={draw,rounded corners=2pt,fill=dxgrey,minimum width=2.1cm,minimum height=0.6cm},
  m/.style={-{Stealth[length=2mm]},thick,dxblue},
  r/.style={-{Stealth[length=2mm]},dashed,gray},
  lb/.style={fill=white,inner sep=1pt,midway,above=1.5pt}]
\node[life] (u) at (0,0) {Người dùng};
\node[life] (p) at (2.6,0) {Portal};
\node[life] (k) at (5.2,0) {Keycloak};
\node[life] (g) at (7.8,0) {APISIX};
\node[life] (c) at (10.4,0) {Core};
\node[life] (d) at (13.0,0) {PostgreSQL};
\foreach \n in {u,p,k,g,c,d} \draw[gray,dotted] (\n.south) -- ++(0,-7.5);
\draw[m] (0,-1.0) -- node[lb]{1. mở portal, đăng nhập} (2.6,-1.0);
\draw[m] (2.6,-1.7) -- node[lb]{2. Authorization Code + PKCE} (5.2,-1.7);
\draw[r] (5.2,-2.3) -- node[lb]{token (giữ phía máy chủ)} (2.6,-2.3);
\draw[m] (2.6,-3.0) -- node[lb]{3. POST /api/workflows + Bearer} (7.8,-3.0);
\node[draw,fill=white,rounded corners=2pt,align=center] at (7.8,-3.75) {4. kiểm JWT (JWKS),\\giới hạn tốc độ};
\draw[m] (7.8,-4.5) -- node[lb]{5. chuyển tiếp} (10.4,-4.5);
\node[draw,fill=white,rounded corners=2pt,align=center] at (10.4,-5.25) {6. iss/aud/exp, 5 lớp quyền,\\JSON Schema, chạy DMN};
\draw[m] (10.4,-6.25) -- node[lb]{7. lưu phiếu, việc, audit} (13.0,-6.25);
\draw[r] (10.4,-7.0) -- node[lb]{8. JSON đã che dữ liệu cá nhân} (2.6,-7.0);
\end{tikzpicture}
\caption{Luồng báo hỏng từ portal tới cơ sở dữ liệu.}
\end{figure}

APISIX chặn sớm, Core quyết định độc lập: gateway cấu hình sai cũng không mở được quyền mà Core không cho. Portal và MCP
chỉ ẩn nút theo vai trò để thuận tiện; quyết định cho phép hay từ chối luôn ở Core.

%================ 3 ================
\section{Không gian H-P-D-I}
Bốn trục của DX-OS gắn với nhau qua một API REST duy nhất, bảo vệ bởi APISIX và Keycloak.
\begin{description}[leftmargin=0pt,labelsep=0.5em]
\item[\textbf{H} Con người] Portal tiếng Việt (Prefab UI qua BFF; token OAuth2 ở phiên phía máy chủ, trình duyệt chỉ giữ cookie \code{httpOnly}),
sáu tab gồm Việc của tôi, Báo hỏng, Tra cứu, Dashboard, Sơ đồ quy trình, Trạng thái hệ thống. Máy chủ MCP (stdio) có năm
công cụ \code{bao_hong}, \code{viec_cua_toi}, \code{tra_cuu_phieu}, \code{tim_huong_dan}, \code{hoan_thanh_viec}; thao tác ghi
bắt buộc hỏi xác nhận. CLI \code{dxlab}: \code{login} (device flow), \code{whoami}, \code{validate}, \code{publish},
\code{logout}. OTP (TOTP) bắt buộc cho vai trò duyệt và quản trị.
\item[\textbf{P} Quy trình] BPMN 2.0 chạy bằng SpiffWorkflow nhúng trong Core; hộp việc \cmd{GET /api/tasks} theo vai trò và phạm vi đơn vị; luật không
kiêm nhiệm SoD-1 đến SoD-5; quản trị thay đổi định nghĩa (mục~\ref{sec:def}).
\item[\textbf{D} Dữ liệu] PostgreSQL là kho duy nhất, chỉ Core kết nối. Form theo JSON Schema; Core từ chối trường lạ ngoài schema. Trường dữ liệu cá nhân
gắn \code{x-personal-data} và bị che khi người xem không có quyền. Nhật ký audit ghi thời điểm, người làm, vai trò, hành động,
căn cứ thẩm quyền.
\item[\textbf{I} Trí tuệ] Bảng luật DMN 1.3 có mã lý do (không hộp đen) và tìm kiếm kết hợp FTS tiếng Việt với embedding cục bộ (mục~\ref{sec:ai}). Quyết
định cuối thuộc người có thẩm quyền.
\end{description}

\subsection{Quản trị thay đổi định nghĩa}\label{sec:def}
Quy trình, luật và form là dữ liệu, nên thay đổi chúng được kiểm soát như thay đổi mã
(\code{core/src/dxlab_core/definitions/router.py}, \code{CHANGELOG.md} mục \code{[0.2.0]}):
\begin{enumerate}
\item \textbf{Nạp}: \cmd{POST /api/definitions/{app_id}/versions} hoặc \cmd{dxlab publish apps/<app>}. Validator từ chối
script, trả \code{DEFINITION_REJECTED}; mỗi lần nạp là một phiên bản mới.
\item \textbf{Phân quyền nạp}: \code{rule_admin} chỉ nạp thay đổi \code{.dmn}; thay đổi khác cần \code{process_admin}.
\item \textbf{Duyệt}: \code{.../versions/{n}/activate} do \code{process_owner} làm và phải là người khác người nạp (SoD-5).
\item \textbf{Áp dụng}: phiếu mới dùng phiên bản mới, phiếu cũ giữ luật của phiên bản đã tạo ra nó.
\end{enumerate}

%================ 4 ================
\section{Ứng dụng nghiệp vụ}
\subsection{lab-repair: sáu bước từ báo hỏng tới nghiệm thu}
\begin{figure}[ht]
\centering
\begin{tikzpicture}[font=\sffamily\scriptsize,
  st/.style={draw,rounded corners=3pt,fill=t3,minimum width=2.7cm,minimum height=1.0cm,align=center,text width=2.6cm,inner sep=2pt},
  dm/.style={draw,rounded corners=3pt,fill=t1,minimum width=2.45cm,align=center,font=\sffamily\tiny},
  a/.style={-{Stealth[length=2mm]},thick,dxblue}]
\node[st] (s1) at (0,0) {1. Báo hỏng\\\emph{report}};
\node[st] (s2) at (3.2,0) {2. Kiểm tra, biên bản\\{\tiny\emph{pending\_assessment}}};
\node[st] (s3) at (6.4,0) {3. Xác nhận nhu cầu\\{\tiny\emph{awaiting\_endorsement}}};
\node[st] (s4) at (9.6,0) {4. Duyệt chi phí\\{\tiny\emph{awaiting\_approval}}};
\node[st] (s5) at (9.6,-2.3) {5. Sửa chữa\\\emph{repairing}};
\node[st] (s6) at (6.4,-2.3) {6. Nghiệm thu\\{\tiny\emph{awaiting\_acceptance}}};
\node[st,fill=t2] (done) at (3.2,-2.3) {Hoàn tất\\\emph{completed}};
\draw[a] (s1)--(s2); \draw[a] (s2)--(s3); \draw[a] (s3)--(s4); \draw[a] (s4)--(s5); \draw[a] (s5)--(s6); \draw[a] (s6)--(done);
\node[dm] (d1) at (0,-1.55) {triage.dmn: P1/P2/P3\\routing.dmn: đơn vị, vai trò};
\node[dm] (d4) at (7.9,1.45) {authority.dmn: ai duyệt, bao nhiêu báo giá\\overrun.dmn: vượt chi phí thì duyệt lại};
\draw[dashed,gray] (s1) -- (d1); \draw[dashed,gray] (s4) -- (d4);
\draw[a,dashed] (s2.south) -- (3.2,-1.0) -- (3.2,-3.6) -- (9.6,-3.6) -- (s5.south);
\node[font=\sffamily\tiny,fill=white] at (6.4,-3.6) {không cần duyệt (chi phí bằng 0): sang sửa chữa};
\end{tikzpicture}
\caption{Sáu bước của \code{lab-repair}, tên trạng thái theo \code{lab-repair.bpmn}. Khi không cần duyệt, phiếu đi thẳng từ kiểm tra sang sửa chữa; nghiệm thu không đạt thì quay lại sửa.}
\end{figure}

Bước 1 kích hoạt \code{triage.dmn} và \code{routing.dmn}, rồi phiếu dừng ở \code{pending_assessment}. Bước 2: kỹ thuật viên
đơn vị được định tuyến kiểm tra và lập biên bản. Kết quả có bốn hướng: sửa không tốn chi phí, sửa có chi phí, không sửa được
(chuyển Ban CSVC lập đề nghị mua thay) hoặc báo nhầm (đóng phiếu). Sau biên bản, \code{authority.dmn} xác định có cần duyệt,
ai duyệt và cần mấy báo giá; thiếu báo giá thì phiếu quay về kiểm tra. Nếu cần duyệt, bước 3 là trưởng phòng lab xác nhận nhu
cầu (không cần sửa thì đóng), bước 4 là người có thẩm quyền duyệt (bảng~\ref{tab:auth}). Bước 5 sửa và ghi chi phí thực tế;
\code{overrun.dmn} yêu cầu duyệt lại nếu chi phí thực tế vượt số đã duyệt quá 10\% hoặc vượt hạn mức người đã duyệt. Bước 6:
người khác người sửa nghiệm thu (SoD-2); không đạt thì phiếu quay lại sửa.

\textbf{Các bảng luật.}
\code{triage.dmn} nhận \code{safety_hazard}, \code{impact}, \code{urgency} và trả \code{priority}, \code{priority_reason},
hạn kiểm tra và hạn hoàn tất: P1 (nguy cơ an toàn, hoặc cả phòng không dùng được) kiểm tra trong 4 giờ, hoàn tất trong 24 giờ;
P2 trong 24/72 giờ; P3 trong 48/120 giờ. Mã lý do gồm \code{SAFETY_RISK}, \code{ROOM_BLOCKED}, \code{SESSION_PARTIAL},
\code{REPORTER_URGENT}, \code{DEFAULT}. \code{routing.dmn} ánh xạ loại thiết bị sang đơn vị: CNTT (\code{computer},
\code{projector}, \code{network}) tới \code{/haui/tt-cntt}; điện, điều hòa, chiếu sáng, nước, nội thất tới
\code{/haui/ban-csvc}; thiết bị chuyên dụng \code{lab_equipment} tới khoa sở hữu phòng; loại khác về Ban CSVC để phân loại
thủ công. Mọi con số SLA là \gd{} của đội.

\begin{table}[ht]
\centering\tbl
\caption{Thẩm quyền duyệt sửa chữa (\code{authority.dmn}, \code{docs/AUTHZ.md} mục 5). Hạn mức và mã căn cứ \code{GĐ-...} là giả định, không phải quy chế thật của HaUI.}\label{tab:auth}
\begin{tabularx}{\textwidth}{@{}L{3.4cm}L{3.6cm}L{2.6cm}Y@{}}
\toprule
Chi phí dự kiến (đồng) & Người duyệt & Báo giá tối thiểu & Hình thức \\
\midrule
0 & không cần duyệt & --- & --- \\
$\le 5\,000\,000$ & \code{faculty_head} (khoa sở hữu phòng) & 1 (hoặc ước tính) & theo phân cấp \\
trên 5 triệu đến $20\,000\,000$ & \code{school_head} (trường thành viên) & 1 (dưới 20 triệu), 3 (đúng 20 triệu) & theo phân cấp \\
trên 20 triệu đến $100\,000\,000$ & \code{vice_president} & 3 & KT. Giám đốc \\
$> 100\,000\,000$ & \code{president} & 3 & trực tiếp \\
mua thay thế, cải tạo & ngoài hệ thống (phiếu đề nghị gửi Ban CSVC) & --- & --- \\
\bottomrule
\end{tabularx}
\end{table}

\subsection{room-booking}
Gói dữ liệu mượn phòng thực hành dùng danh mục phòng của \code{lab-repair} qua \code{record_bindings}. Người đề nghị khai báo
\code{schedule_conflict}; ứng dụng \emph{không} tự kiểm tra lịch hay chặn đặt trùng. Quyết định duyệt hoặc từ chối đều bắt buộc
ghi lý do. Điểm cần chứng minh là ứng dụng này nạp bằng \cmd{dxlab publish} mà không đổi một dòng mã Core.

%================ 5 ================
\section{Phân quyền năm lớp và ủy quyền TUQ.}
Core kiểm theo thứ tự, mặc định từ chối (chi tiết, căn cứ và nhãn ``giả định'' ở \code{docs/AUTHZ.md}).
\begin{center}\tbl
\begin{tabularx}{\textwidth}{@{}L{2.6cm}L{4.3cm}L{3.9cm}Y@{}}
\toprule
Lớp & Câu hỏi & Nguồn dữ liệu & Mã lỗi \\
\midrule
1. Ai là ai & Token hợp lệ không, vai trò gì ở đơn vị nào? & nhóm Keycloak & \code{UNAUTHENTICATED}, \code{FORBIDDEN_ROLE} \\
2. Phạm vi & Đối tượng thuộc đơn vị của người làm (kể cả đơn vị con)? & cây đơn vị trong Core & \code{FORBIDDEN_SCOPE} \\
3. Thẩm quyền & Hạn mức đủ không, hoặc có ủy quyền còn hạn? & \code{authority.dmn}, \code{/api/delegations} & \code{FORBIDDEN_AUTHORITY} \\
4. Không kiêm nhiệm & Có đang tự duyệt hay tự nghiệm thu việc mình làm? & \code{dxlabExcludeActorsOf} trên bước BPMN & \code{FORBIDDEN_SOD} \\
5. Dữ liệu cá nhân & Có được thấy liên hệ người báo? & \code{x-personal-data} trong schema & trường bị che \\
\bottomrule
\end{tabularx}
\end{center}
\textbf{TUQ.} (\code{/api/delegations}, issue \#18): người giữ vai trò duyệt giao cho người cùng đơn vị duyệt thay tối đa
30 ngày, bắt buộc số giấy ủy quyền, thu hồi được bất kỳ lúc nào; không vượt vai trò và hạn mức của người ủy quyền; người nhận
không ủy quyền lại được; phiếu và audit ghi hình thức TUQ., người ủy quyền và số giấy. Kịch bản kiểm thử số 7 của
\code{docs/AUTHZ.md} mục 12: \code{ptk.hoa} duyệt được 3.500.000 đồng theo TUQ. của \code{tk.hai}, duyệt 12.000.000 đồng bị
từ chối vì vượt hạn mức người ủy quyền, và sau hạn cũng bị từ chối. Mười bốn kịch bản phân quyền này chạy trên gói
\code{lab-repair} thật. \emph{Hiện trạng}: chưa có màn hình tạo hay thu hồi ủy quyền trên portal; dùng API
(\code{/api/docs}).

\section{Luồng dữ liệu}
\begin{figure}[ht]
\centering
\begin{tikzpicture}[font=\sffamily\scriptsize,
  p/.style={draw,rounded corners=3pt,fill=t1,minimum width=2.6cm,minimum height=0.9cm,align=center},
  s/.style={draw,fill=t3,minimum width=2.6cm,minimum height=0.9cm,align=center},
  a/.style={-{Stealth[length=2mm]},thick,dxblue}]
\node[p] (ng) at (0,0) {Người dùng\\form + Bearer JWT};
\node[p] (gw) at (3.9,0) {APISIX\\kiểm JWT, giới hạn tốc độ};
\node[p] (v) at (7.8,0) {Core: kiểm schema\\và 5 lớp quyền};
\node[p] (w) at (7.8,-1.8) {SpiffWorkflow\\BPMN + DMN};
\node[s] (pg) at (3.9,-1.8) {PostgreSQL\\phiếu, việc, audit, FTS};
\node[s] (zv) at (3.9,-3.4) {Chỉ mục vector zvec\\384 chiều (volume)};
\draw[a] (ng)--(gw); \draw[a] (gw)--(v); \draw[a] (v)--(w); \draw[a] (w)--(pg);
\draw[a] (pg)--(zv);
\draw[a,gray,dashed] (v.north) to[bend right=18] node[above,font=\sffamily\tiny]{JSON đã che dữ liệu cá nhân} (ng.north);
\end{tikzpicture}
\caption{Đường đi của dữ liệu một phiếu. Chỉ mục vector chỉ dùng cho tìm kiếm SOP.}
\end{figure}
Dữ liệu vào Core phải khớp JSON Schema của bước; trường lạ bị từ chối nên client không ghi được biến nội bộ của quy trình.
Trạng thái quy trình, phiếu, việc và audit nằm trong PostgreSQL; tài liệu SOP và chỉ mục vector nằm trên volume do Core quản lý.
Phản hồi cho người xem luôn đi qua bước che dữ liệu cá nhân theo quyền của người đó (Nghị định 13/2023/NĐ-CP là căn cứ
thiết kế; \code{docs/AUTHZ.md}).

%================ 7 AI ================
\section{Trí tuệ nhân tạo và học máy}\label{sec:ai}
\textbf{Nguyên tắc}: quyết định liên quan an toàn, hạn mức và nghiệm thu dùng luật DMN có mã lý do, không dùng mô hình hộp
đen; phần học máy chỉ gợi ý tri thức. Quyền quyết định luôn thuộc người có thẩm quyền.

\textbf{Mô hình}: \code{sentence-transformers/paraphrase-multilingual-MiniLM-L12-v2} (Apache-2.0, ONNX, 384 chiều, khoảng
470 MB), nạp offline bằng \code{scripts/download-models.sh}, chạy qua \code{fastembed}, lưu vector bằng thư viện \code{zvec}.
Không dữ liệu nào gửi ra ngoài. Bỏ bước tải mô hình thì \code{/api/search} tự lùi về toàn văn FTS.

\textbf{Xếp hạng kết hợp}: điểm $= 0{,}7\cdot\dfrac{r}{r+0{,}1} + 0{,}3\cdot\max\!\bigl(0,\,1-\dfrac{d}{0{,}8}\bigr)$,
với $r$ là \code{ts_rank_cd} của PostgreSQL FTS (có bỏ dấu tiếng Việt) và $d$ là khoảng cách cosine; đoạn có $d>0{,}8$ bị loại
để câu hỏi ngoài phạm vi trả rỗng.

\textbf{Dữ liệu đánh giá} (\code{apps/lab-repair/eval/}, \code{docs/AI.md}): 60 phiếu \emph{tổng hợp} gán nhãn tay (P1 11,
P2 21, P3 28) và 20 câu hỏi tìm kiếm trên 8 SOP của \code{data/sop/} (18 câu có đáp án, 2 câu ngoài phạm vi).

\begin{table}[ht]
\centering\tbl
\caption{Số đo thật (\code{docs/AI.md}); lệnh tái lập ở phụ lục~\ref{sec:chay}.}
\begin{tabular}{@{}lrr@{}}
\toprule
Chỉ số & Đúng/Tổng & Tỷ lệ \\
\midrule
Mức ưu tiên (\code{triage.dmn}) & 53/60 & 88,3\% \\
Đơn vị kỹ thuật (\code{routing.dmn}, \code{tech_unit}) & 57/60 & 95,0\% \\
Vai trò xử lý (\code{tech_role}) & 50/60 & 83,3\% \\
Định tuyến đúng cả đơn vị và vai trò & 50/60 & 83,3\% \\
Tìm kiếm TF-IDF, top-1 (18 câu có đáp án) & 17/18 & 94,4\% \\
Tìm kiếm kết hợp \code{/api/search}, top-1 và top-3 & 18/18 & 100\% \\
Câu ngoài phạm vi trả rỗng: TF-IDF / kết hợp & 0/2 / 2/2 & 0\% / 100\% \\
\bottomrule
\end{tabular}
\end{table}

Ma trận nhầm lẫn mức ưu tiên (hàng là nhãn tay, cột là DMN): P1 $\to$ (8, 1, 2); P2 $\to$ (1, 19, 1); P3 $\to$ (0, 2, 26).
Câu cải thiện rõ nhất là ``wifi chập chờn rớt liên tục'': TF-IDF xếp SOP đúng ở hạng 7 vì từ ``chập chờn'' kéo sang SOP chiếu
sáng, tìm kiếm kết hợp xếp hạng 1.

\textbf{Giới hạn} (nêu trước hết ở \code{docs/AI.md}):
\begin{itemize}
\item Nhãn do một agent soạn hộ, \emph{chưa có người rà}; phiếu là dữ liệu tổng hợp, không phải phiếu thật của HaUI.
\item Nhãn định tuyến và \code{routing.dmn} cùng nguồn (\code{docs/AUTHZ.md} mục 6), nên 95\% chủ yếu đo \emph{đầu vào người
báo chọn} chứ không đo độ ``thông minh'' của luật. 60 phiếu là mẫu nhỏ: một phiếu sai đổi 1,7 điểm phần trăm.
\item Sai nghiêm trọng nhất (T03, T07, T10): người báo không đánh dấu nguy cơ an toàn dù mô tả có khói hoặc lộ dây điện, DMN
không đọc được văn bản tự do nên xếp dưới mức; cần form bắt buộc xác nhận an toàn hoặc kỹ thuật viên nâng mức ở bước kiểm
tra, chưa phải việc \code{triage.dmn} tự giải được.
\item 20 câu hỏi do người đã biết nội dung SOP viết, nên mốc so sánh có thể lạc quan; chưa có câu do người ngoài dự án soạn.
\end{itemize}

%================ 8 LICENSE ================
\section{Giấy phép và phụ thuộc}
\begin{center}\tbl
\begin{tabularx}{\textwidth}{@{}L{4.3cm}L{3.5cm}Y@{}}
\toprule
Tài sản & Giấy phép & Cách khai báo \\
\midrule
Mã nguồn, script, cấu hình build & Apache-2.0 & header SPDX trong tệp \\
Tài liệu (README, CHANGELOG, báo cáo) & CC BY 4.0 & header SPDX trong tệp \\
Tệp không chứa được comment (\code{.bpmn}, \code{.dmn}, \code{.json}, \code{uv.lock}, PDF) & theo loại tệp & chú thích tập trung trong \code{REUSE.toml} \\
Dữ liệu bên thứ ba & giữ giấy phép gốc & \code{DATA_SOURCES.yml} \\
\bottomrule
\end{tabularx}
\end{center}
Repo tuân thủ REUSE 3.3 (\cmd{reuse lint} xanh, và là một job CI bắt buộc). Bảng đầy đủ ở \code{docs/DEPENDENCIES.md}: mục 2
là năm image container ghim theo digest \code{sha256} (PostgreSQL, Keycloak, APISIX, image nền Python 3.12, \code{uv}); mục 3 là
mô hình embedding (Apache-2.0); mục 4 là khai báo thành phần đóng gói sẵn theo quy ước Fedora cho bộ giao diện của
\code{prefab-ui==0.21.1} (Apache-2.0); mục 5 là toàn bộ gói Python lấy từ \code{uv.lock}, có phiên bản và giấy phép. SpiffWorkflow
(LGPL-3.0) được dùng như thư viện, không sửa mã. Không commit thư viện chép tay. Cấu hình mẫu ở \code{.env.example};
\code{.env} không commit, \code{scripts/up.sh} sinh ngẫu nhiên các secret để trống và job gitleaks quét secret trên mỗi PR.

%================ 9 BẰNG CHỨNG ================
\section{Bằng chứng nguồn mở}
Số liệu lấy ngày 01/11/2026 từ \code{git} và \code{gh}; các số tăng theo thời gian.
\begin{center}\tbl
\begin{tabularx}{\textwidth}{@{}L{3.3cm}Y@{}}
\toprule
Bằng chứng & Số liệu, vị trí kiểm \\
\midrule
Repo công khai, lịch sử thật & 115 commit tới tag \code{v0.2.1}; commit đầu \code{5d15b80} ngày 27/10/2026 (\cmd{git log}) \\
Release có phiên bản & \code{v0.2.0} (30/10/2026) và \textbf{\code{v0.2.1}} (01/11/2026, giờ Việt Nam). Tệp kèm \code{v0.2.1}:
\code{dx-lab-v0.2.1.tar.gz}, \code{SHA256SUMS}, SBOM CycloneDX \code{dx-lab-v0.2.1.cdx.json} và SPDX \code{dx-lab-v0.2.1.spdx.json}
(syft 1.55.0); kiểm bằng \cmd{sha256sum -c SHA256SUMS}; \cmd{gh release view v0.2.1} \\
Quy trình phát hành & \code{.github/workflows/release.yml}; \code{CHANGELOG.md} theo Keep a Changelog \\
CI bắt buộc & Năm job trong \code{ci.yml}: REUSE; build wheel/sdist, cài sạch và import ngoài cây mã; DCO
(\code{Signed-off-by} từng commit); gitleaks; pytest toàn workspace cùng \cmd{dxlab validate apps/} \\
Issue, PR, review & 43 PR đã gộp, 57 issue (mọi trạng thái, đếm bằng \code{gh} ngày 01/11/2026). Mỗi thay đổi gắn issue, PR theo
mẫu, review của người khác tác giả, chỉ maintainer gộp bằng rebase (\code{CONTRIBUTING.md}) \\
Tái lập từ source & \cmd{uv sync --all-packages}, \cmd{uv build --all-packages}, \code{scripts/up.sh} (\code{README.md}) \\
QA-01 (\code{docs/qa/QA-01.md}) & Dựng từ \cmd{git clone} theo README rồi chạy \code{demo-e2e.sh}: lần 1 lỗi ở bước 2 (script đọc sai tên trường,
issue \#118, sửa ở PR \#119); lần 2 đạt, phiếu tới \code{completed}. \textbf{Chạy trên máy của mentor} (macOS 27.2, Apple Silicon,
Docker 29.8.2), sau khi xóa container, volume, image và build cache của dự án; \emph{chưa} chạy trên máy khác hay Linux/WSL2 \\
\bottomrule
\end{tabularx}
\end{center}

\section{Đánh giá và hạn chế}
\textbf{Đã đạt.} Core headless tách khỏi giao diện và ứng dụng; ứng dụng thứ hai nạp không sửa Core; phân quyền năm lớp kèm
TUQ. được kiểm bằng 14 kịch bản trên gói thật; vòng đời phiếu chạy đầu-cuối qua APISIX với Keycloak thật (\code{scripts/demo-e2e.sh});
repo tuân thủ REUSE, có release, SBOM và checksum.

\textbf{Hạn chế đã biết} (\code{CHANGELOG.md} mục \code{[0.2.1]}):
\begin{itemize}
\item Hạn mức, cơ cấu tổ chức, SLA, tài khoản demo và các căn cứ \code{GĐ-...} là \emph{giả định} của đội, không phải quy chế thật của HaUI.
\item Dữ liệu đánh giá AI là tổng hợp, nhãn chưa có người rà (mục~\ref{sec:ai}).
\item Chỉ mục toàn văn lấy 150.000 ký tự đầu của tài liệu.
\item Docker Compose là bản demo một máy; Keycloak chạy \code{start-dev}; triển khai thật cần TLS và cơ sở dữ liệu riêng
(\code{docs/ARCHITECTURE.md} mục 4).
\item Portal chưa có màn hình tạo hoặc thu hồi ủy quyền.
\item QA-01 chỉ chạy trên một máy (máy mentor).
\end{itemize}

\section{Lộ trình}\label{sec:lotrinh}
Các mục sau \emph{chưa có} trong \code{v0.2.1}, không được coi là tính năng đã chạy; Core không gọi LLM theo thiết kế.
\begin{itemize}
\item Chạy MCP với host và LLM mở (issue \#36).
\item Che dữ liệu cá nhân trong văn bản tự do bằng mô hình NLP OpenMed (issue \#97).
\item Đo và tối ưu migration \code{0006} trên bảng \code{documents} lớn (issue \#96).
\item Màn hình tạo và thu hồi ủy quyền trên portal; rà nhãn bộ 60 phiếu bởi người am hiểu phòng lab; QA trên máy thứ hai.
\end{itemize}

%================ PHỤ LỤC ================
\appendix
\section{Chạy nhanh}\label{sec:chay}
Yêu cầu: Linux, macOS hoặc Windows (WSL2); Docker với Compose v2; \code{git}, \code{curl}, \code{openssl}, \code{bash}.
\begin{Verbatim}[fontsize=\small,frame=single,framesep=2mm,rulecolor=\color{gray}]
git clone https://github.com/HaUI-NPC/dx-lab.git && cd dx-lab
scripts/up.sh
scripts/download-models.sh
scripts/seed.sh
\end{Verbatim}
\code{up.sh} tạo \code{.env}, build, chạy 5 dịch vụ và chờ healthy; \code{download-models.sh} (tùy chọn) tải mô hình embedding; \code{seed.sh} nạp \code{lab-repair}, danh mục, SOP và phiếu mẫu. Sau đó: portal \url{http://localhost:8000}, Swagger \url{http://localhost:9080/api/docs}, Keycloak \url{http://localhost:8180}.
Kiểm đầu-cuối bằng \code{scripts/demo-e2e.sh}; đặt lại bằng \cmd{scripts/reset.sh --yes}. Tái lập số đo AI:
\begin{Verbatim}[fontsize=\small,frame=single,framesep=2mm,rulecolor=\color{gray}]
uv run --with "spiffworkflow>=3.2" python apps/lab-repair/eval/measure_dmn.py
python3 apps/lab-repair/eval/measure_search_baseline.py --sop-dir data/sop
\end{Verbatim}

\section{Tài khoản demo (giả định)}
Toàn bộ tài khoản, tên và hạn mức dưới đây là \textbf{giả định} của đội phục vụ chạy thử, không phải cán bộ thật của HaUI.
README có đủ 17 tài khoản; mật khẩu chung nằm ở \code{DXLAB_DEMO_PASSWORD} trong \code{.env}, khóa OTP ở
\code{DXLAB_DEMO_OTP_SECRET}.
\begin{center}\tbl
\begin{tabularx}{\textwidth}{@{}L{2.2cm}L{4.8cm}Y L{1.1cm}@{}}
\toprule
Tài khoản & Vai trò & Việc dùng trong kịch bản & OTP \\
\midrule
\code{sv.an} & \code{reporter} (sinh viên) & báo hỏng & không \\
\code{ktv.dung} & \code{facility_technician} & lập biên bản, sửa chữa & không \\
\code{tl.giang} & \code{lab_manager} & xác nhận nhu cầu, nghiệm thu & không \\
\code{tk.hai} & \code{faculty_head} & duyệt đến 5.000.000 đ & có \\
\code{ptk.hoa} & \code{faculty_deputy} & nhận ủy quyền TUQ. & có \\
\code{ht.khanh} & \code{school_head} & duyệt đến 20.000.000 đ & có \\
\code{pgd.minh} & \code{vice_president}, \code{process_owner} & duyệt đến 100.000.000 đ, duyệt phiên bản & có \\
\code{qt.son} & \code{process_admin} & nạp ứng dụng qua CLI & có \\
\code{admin.dx} & \code{system_admin} & không làm được việc nghiệp vụ (SoD-4) & có \\
\bottomrule
\end{tabularx}
\end{center}

\end{document}
